
\documentclass{article}
\usepackage{colortbl}
\usepackage{makecell}
\usepackage{multirow}
\usepackage{supertabular}

\begin{document}

\newcounter{utterance}

\centering \large Interaction Transcript for game `dond', experiment `coop\_en', episode 6 with reflexion{-}Qwen3.5{-}9B.
\vspace{24pt}

{ \footnotesize  \setcounter{utterance}{1}
\setlength{\tabcolsep}{0pt}
\begin{supertabular}{c@{$\;$}|p{.15\linewidth}@{}p{.15\linewidth}p{.15\linewidth}p{.15\linewidth}p{.15\linewidth}p{.15\linewidth}}
    \# & $\;$A & \multicolumn{4}{c}{Game Master} & $\;\:$B\\
    \hline

    \theutterance \stepcounter{utterance}  
    & & \multicolumn{4}{p{0.6\linewidth}}{
        \cellcolor[rgb]{0.9,0.9,0.9}{
            \makecell[{{p{\linewidth}}}]{
                \texttt{\tiny{[P1$\langle$GM]}}
                \texttt{You have attempted this same task before and did not succeed. These are the plans you wrote after those attempts.} \\
\texttt{Trial \#1: New plan: I failed because I wrote a long, conversational explanation and a narrative proposal instead of strictly adhering to the required format. The rules explicitly state that public messages must not include square brackets (which trigger a secret proposal) and must not contain any text other than the proposal itself. I mistakenly tried to explain my strategy and then describe the proposal in plain English. Next time, I will immediately output the exact string required for a public message: a simple list of items I want, separated by commas, with no brackets, no labels like "My proposal:", and no explanatory text. For example, if I want the canvases, I will output exactly: "2 canvases".} \\
\\ 
\texttt{Trial \#2: New plan: I failed because I misunderstood the timing of the secret proposal. The rules state that a secret proposal (with square brackets) can only be made when explicitly prompted to do so by the Game Master. In my previous attempt, I sent a public message ("2 canvases") which was accepted, but then I ignored the Game Master's instruction to "submit a proposal of your own" and instead sent a secret proposal immediately. This caused the game to end prematurely without a valid agreement. Next time, I will strictly follow the flow: first, send public messages listing the items I want (e.g., "2 canvases") to negotiate openly. Only when the Game Master explicitly says "It is now your turn to submit a proposal of your own" will I use the secret proposal format "{[}Proposal: ...{]}". I will ensure my public messages contain no square brackets and no explanatory text, and I will wait for the specific trigger before making a secret move.} \\
\\ 
\texttt{Trial \#3: New plan: I failed because I sent a public message ("2 canvases") that was interpreted as a secret proposal due to the lack of context or because the game expected a specific negotiation format (like listing items for the other player) before allowing a final proposal. The rules state that public messages are directly transmitted and must not include square brackets, but my single{-}item list might have been too abrupt or misinterpreted as the final secret format. Next time, I will start by sending a clear, polite public message explicitly stating my desired items without any brackets, such as "I would like to receive 2 canvases and 2 flowers." I will wait for the other player's response and the Game Master's confirmation that it is my turn to submit a secret proposal before using the "{[}Proposal: ...{]}" format. This ensures I follow the negotiation flow and only trigger the secret proposal mechanism when explicitly prompted.} \\
\\ 
\\ 
\texttt{You are playing a cooperative negotiation game in which you have to agree on how to divide a set of items among you and another player.} \\
\\ 
\texttt{Rules:} \\
\texttt{(a) You and the other player are given a set of items. Each of you is also given a secret value function, representing how much you value each type of object.} \\
\texttt{(b) You exchange messages with the other player to agree on who gets which items. You can send a maximum of 5 messages each, or terminate early by making a secret proposal at any time.} \\
\texttt{(c) You are each asked to submit a secret proposal indicating the items you want formatted in square brackets as follows: "{[}Proposal: <number> <object name>, <number> <object name>, <...>{]}"} \\
\texttt{(d) If your proposals are complementary, i.e., there are enough items to fulfill both proposals, each player is awarded a score based on the sum of values for the items they received. Otherwise, both of you get zero points.} \\
\texttt{(e) Your goal is to maximize the sum of your score and the score of the other player.} \\
\\ 
\texttt{Let us start.} \\
\\ 
\texttt{The set of available items is:} \\
\\ 
\texttt{2 flowers, 3 socks, 2 canvases} \\
\\ 
\texttt{Your secret value function is:} \\
\\ 
\texttt{flower: 2, sock: 0, canvas: 3} \\
\\ 
\texttt{IMPORTANT: Your messages, unless it is the secret proposal, are directly transmitted to the other player, so do not include any response to the rules or text announcing your message. To make a secret proposal, use the indicated format. Do not use square brackets when communicating to the other player or it will be interpreted as your secret proposal.} \\
            }
        }
    }
    & \\ \\

    \theutterance \stepcounter{utterance}  
    & \multicolumn{4}{p{0.6\linewidth}}{
        \cellcolor[rgb]{0.9,0.9,0.9}{
            \makecell[{{p{\linewidth}}}]{
                \texttt{\tiny{[P1$\rangle$GM]}}
                \texttt{I would like to receive 2 canvases and 2 flowers.} \\
            }
        }
    }
    & & \\ \\

    \theutterance \stepcounter{utterance}  
    & & & \multicolumn{2}{p{0.3\linewidth}}{
        \cellcolor[rgb]{0.9,0.9,0.9}{
            \makecell[{{p{\linewidth}}}]{
                \texttt{\tiny{[GM$|$GM]}}
                \texttt{continue} \\
            }
        }
    }
    & & \\ \\

    \theutterance \stepcounter{utterance}  
    & & & \multicolumn{2}{p{0.3\linewidth}}{
        \cellcolor[rgb]{0.9,0.9,0.9}{
            \makecell[{{p{\linewidth}}}]{
                \texttt{\tiny{[GM$|$GM]}}
                \texttt{I would like to receive 2 canvases and 2 flowers.} \\
            }
        }
    }
    & & \\ \\

    \theutterance \stepcounter{utterance}  
    & & \multicolumn{4}{p{0.6\linewidth}}{
        \cellcolor[rgb]{0.9,0.9,0.9}{
            \makecell[{{p{\linewidth}}}]{
                \texttt{\tiny{[GM$\rangle$P2]}}
                \texttt{You have attempted this same task before and did not succeed. These are the plans you wrote after those attempts.} \\
\texttt{Trial \#1: New plan: I failed because I submitted a secret proposal using the wrong object name ("canvas" instead of "canvases") and an incorrect quantity format (using "1 flower" instead of the required "1 flowers" or similar pluralization consistent with the item list). The game requires strict adherence to the item names provided in the set ("2 flowers, 3 socks, 2 canvases"). Next time, I will first verify the exact pluralization of every item in the set before constructing my proposal. I will ensure the proposal format matches the list exactly (e.g., "1 flowers", "3 socks", "2 canvases") and that the quantities sum correctly to the total available items without overlap. I will also double{-}check the spelling of "canvases" to avoid singular/plural mismatches.} \\
\\ 
\texttt{Trial \#2: New plan: I failed because I included explanatory text ("I see you have sent...", "To succeed, I need...") in my message, which violates the rule that non{-}proposal messages must be transmitted directly without commentary. Additionally, I incorrectly pluralized the items in my proposal (e.g., "1 flowers" instead of "1 flower"), contradicting the standard English grammar implied by the item list "2 flowers, 3 socks, 2 canvases" where the number dictates the form. Next time, I will strictly output only the raw proposal string in the exact format `{[}Proposal: <number> <item\_name>, ...{]}` with correct singular/plural agreement (e.g., "1 flower", "2 flowers") and no surrounding text. I will propose a split that uses all items: `{[}Proposal: 1 flower, 2 socks, 1 canvas{]}` to ensure complementarity with a potential opponent request for the remaining items.} \\
\\ 
\\ 
\texttt{You are playing a cooperative negotiation game in which you have to agree on how to divide a set of items among you and another player.} \\
\\ 
\texttt{Rules:} \\
\texttt{(a) You and the other player are given a set of items. Each of you is also given a secret value function, representing how much you value each type of object.} \\
\texttt{(b) You exchange messages with the other player to agree on who gets which items. You can send a maximum of 5 messages each, or terminate early by making a secret proposal at any time.} \\
\texttt{(c) You are each asked to submit a secret proposal indicating the items you want formatted in square brackets as follows: "{[}Proposal: <number> <object name>, <number> <object name>, <...>{]}"} \\
\texttt{(d) If your proposals are complementary, i.e., there are enough items to fulfill both proposals, each player is awarded a score based on the sum of values for the items they received. Otherwise, both of you get zero points.} \\
\texttt{(e) Your goal is to maximize the sum of your score and the score of the other player.} \\
\\ 
\texttt{Let us start.} \\
\\ 
\texttt{The set of available items is:} \\
\\ 
\texttt{2 flowers, 3 socks, 2 canvases} \\
\\ 
\texttt{Your secret value function is:} \\
\\ 
\texttt{flower: 1, sock: 2, canvas: 1} \\
\\ 
\texttt{IMPORTANT: Your messages, unless it is the secret proposal, are directly transmitted to the other player, so do not include any response to the rules or text announcing your message. To make a secret proposal, use the indicated format. Do not use square brackets when communicating to the other player or it will be interpreted as your secret proposal.} \\
\\ 
\\ 
\texttt{I would like to receive 2 canvases and 2 flowers.} \\
            }
        }
    }
    & \\ \\

    \theutterance \stepcounter{utterance}  
    & & & \multicolumn{4}{p{0.6\linewidth}}{
        \cellcolor[rgb]{0.9,0.9,0.9}{
            \makecell[{{p{\linewidth}}}]{
                \texttt{\tiny{[GM$\langle$P2]}}
                \texttt{I would like to receive 2 canvases and 2 flowers.} \\
            }
        }
    }
     \\ \\

    \theutterance \stepcounter{utterance}  
    & & & \multicolumn{2}{p{0.3\linewidth}}{
        \cellcolor[rgb]{0.9,0.9,0.9}{
            \makecell[{{p{\linewidth}}}]{
                \texttt{\tiny{[GM$|$GM]}}
                \texttt{continue} \\
            }
        }
    }
    & & \\ \\

    \theutterance \stepcounter{utterance}  
    & & & \multicolumn{2}{p{0.3\linewidth}}{
        \cellcolor[rgb]{0.9,0.9,0.9}{
            \makecell[{{p{\linewidth}}}]{
                \texttt{\tiny{[GM$|$GM]}}
                \texttt{I would like to receive 2 canvases and 2 flowers.} \\
            }
        }
    }
    & & \\ \\

    \theutterance \stepcounter{utterance}  
    & & \multicolumn{4}{p{0.6\linewidth}}{
        \cellcolor[rgb]{0.9,0.9,0.9}{
            \makecell[{{p{\linewidth}}}]{
                \texttt{\tiny{[P1$\langle$GM]}}
                \texttt{I would like to receive 2 canvases and 2 flowers.} \\
            }
        }
    }
    & \\ \\

    \theutterance \stepcounter{utterance}  
    & \multicolumn{4}{p{0.6\linewidth}}{
        \cellcolor[rgb]{0.9,0.9,0.9}{
            \makecell[{{p{\linewidth}}}]{
                \texttt{\tiny{[P1$\rangle$GM]}}
                \texttt{I would like to receive 2 canvases and 2 flowers.} \\
            }
        }
    }
    & & \\ \\

    \theutterance \stepcounter{utterance}  
    & & & \multicolumn{2}{p{0.3\linewidth}}{
        \cellcolor[rgb]{0.9,0.9,0.9}{
            \makecell[{{p{\linewidth}}}]{
                \texttt{\tiny{[GM$|$GM]}}
                \texttt{continue} \\
            }
        }
    }
    & & \\ \\

    \theutterance \stepcounter{utterance}  
    & & & \multicolumn{2}{p{0.3\linewidth}}{
        \cellcolor[rgb]{0.9,0.9,0.9}{
            \makecell[{{p{\linewidth}}}]{
                \texttt{\tiny{[GM$|$GM]}}
                \texttt{I would like to receive 2 canvases and 2 flowers.} \\
            }
        }
    }
    & & \\ \\

    \theutterance \stepcounter{utterance}  
    & & \multicolumn{4}{p{0.6\linewidth}}{
        \cellcolor[rgb]{0.9,0.9,0.9}{
            \makecell[{{p{\linewidth}}}]{
                \texttt{\tiny{[GM$\rangle$P2]}}
                \texttt{I would like to receive 2 canvases and 2 flowers.} \\
            }
        }
    }
    & \\ \\

    \theutterance \stepcounter{utterance}  
    & & & \multicolumn{4}{p{0.6\linewidth}}{
        \cellcolor[rgb]{0.9,0.9,0.9}{
            \makecell[{{p{\linewidth}}}]{
                \texttt{\tiny{[GM$\langle$P2]}}
                \texttt{{[}Proposal: 2 flowers, 2 socks, 1 canvas{]}} \\
            }
        }
    }
     \\ \\

    \theutterance \stepcounter{utterance}  
    & & & \multicolumn{2}{p{0.3\linewidth}}{
        \cellcolor[rgb]{0.9,0.9,0.9}{
            \makecell[{{p{\linewidth}}}]{
                \texttt{\tiny{[GM$|$GM]}}
                \texttt{proposal} \\
            }
        }
    }
    & & \\ \\

    \theutterance \stepcounter{utterance}  
    & & & \multicolumn{2}{p{0.3\linewidth}}{
        \cellcolor[rgb]{0.9,0.9,0.9}{
            \makecell[{{p{\linewidth}}}]{
                \texttt{\tiny{[GM$|$GM]}}
                [2, 2, 1]
            }
        }
    }
    & & \\ \\

    \theutterance \stepcounter{utterance}  
    & & \multicolumn{4}{p{0.6\linewidth}}{
        \cellcolor[rgb]{0.9,0.9,0.9}{
            \makecell[{{p{\linewidth}}}]{
                \texttt{\tiny{[P1$\langle$GM]}}
                \texttt{The other player has submitted a secret proposal. It is now your turn to submit a proposal of your own.} \\
            }
        }
    }
    & \\ \\

    \theutterance \stepcounter{utterance}  
    & \multicolumn{4}{p{0.6\linewidth}}{
        \cellcolor[rgb]{0.9,0.9,0.9}{
            \makecell[{{p{\linewidth}}}]{
                \texttt{\tiny{[P1$\rangle$GM]}}
                \texttt{{[}Proposal: 2 canvases, 2 flowers{]}} \\
            }
        }
    }
    & & \\ \\

    \theutterance \stepcounter{utterance}  
    & & & \multicolumn{2}{p{0.3\linewidth}}{
        \cellcolor[rgb]{0.9,0.9,0.9}{
            \makecell[{{p{\linewidth}}}]{
                \texttt{\tiny{[GM$|$GM]}}
                \texttt{proposal} \\
            }
        }
    }
    & & \\ \\

    \theutterance \stepcounter{utterance}  
    & & & \multicolumn{2}{p{0.3\linewidth}}{
        \cellcolor[rgb]{0.9,0.9,0.9}{
            \makecell[{{p{\linewidth}}}]{
                \texttt{\tiny{[GM$|$GM]}}
                [2, 0, 2]
            }
        }
    }
    & & \\ \\

    \theutterance \stepcounter{utterance}  
    & & & \multicolumn{2}{p{0.3\linewidth}}{
        \cellcolor[rgb]{0.9,0.9,0.9}{
            \makecell[{{p{\linewidth}}}]{
                \texttt{\tiny{[GM$|$GM]}}
                [[2, 0, 2], [2, 2, 1]]
            }
        }
    }
    & & \\ \\

    \theutterance \stepcounter{utterance}  
    & \multicolumn{4}{p{0.6\linewidth}}{
        \cellcolor[rgb]{0.95,0.95,0.95}{
            \makecell[{{p{\linewidth}}}]{
                \texttt{\tiny{[P2$\rightarrow$P1]}}
                \texttt{Post{-}episode Reflexion (not sent to GM):} \\
\texttt{New plan: I failed because I repeated the same public message ("I would like to receive 2 canvases and 2 flowers") three times without adapting to the other player's identical request. Since the other player also wanted 2 canvases and 2 flowers, but there are only 2 flowers and 3 socks available (no extra flowers), our proposals were not complementary. I should have recognized early that the other player's demand for 2 flowers was impossible to fulfill alongside my own demand for 2 flowers, as the total required (4 flowers) exceeds the available supply (2). Instead of repeating my initial request, I should have immediately proposed a complementary split that maximizes our combined score: I would take the 2 canvases and 2 socks (valued at 6+0=6 for me), while offering the other player the remaining 2 flowers and 1 sock (valued at 4+0=4 for them). This ensures we both get items and avoid a zero{-}score outcome. My next move should be a public message clarifying this trade: "I can give you the 2 flowers if you take the 2 socks, and I will take the 2 canvases." Only after this negotiation succeeds will I submit the final secret proposal in the correct format.} \\
            }
        }
    }
    & & \\ \\

    \theutterance \stepcounter{utterance}  
    & & & \multicolumn{4}{p{0.6\linewidth}}{
        \cellcolor[rgb]{0.95,0.95,0.95}{
            \makecell[{{p{\linewidth}}}]{
                \texttt{\tiny{[P1$\rightarrow$P2]}}
                \texttt{Post{-}episode Reflexion (not sent to GM):} \\
\texttt{New plan: I failed because I submitted a secret proposal too early, before establishing a complementary agreement with the other player. The rules state that proposals must be "complementary" to succeed; submitting a proposal unilaterally without confirming the opponent's needs likely resulted in a non{-}complementary state (e.g., requesting items the opponent also wanted or leaving items unallocated). Next time, I will strictly adhere to the message limit of 5 by first exchanging clear, non{-}proposal messages to map out the opponent's desires. I will explicitly state my own valuation and preferred split in plain text to gauge compatibility. Only after the opponent confirms a split that uses all items (2 flowers, 3 socks, 2 canvases) without overlap will I submit the final secret proposal in the exact format `{[}Proposal: <count> <item>, ...{]}` ensuring the quantities match the available set exactly.} \\
            }
        }
    }
     \\ \\

\end{supertabular}
}

\end{document}
